\documentclass[12pt,twoside,letterpaper]{article}

%NOTE: This report format is 

\newcommand{\reporttitle}{Proyect 1}
\newcommand{\reportauthorOne}{Student 1}
\newcommand{\cidOne}{your id number}
\newcommand{\reportauthorTwo}{Student 2}
\newcommand{\cidTwo}{your id number}
\newcommand{\reporttype}{Coursework}
\bibliographystyle{plain}

% include files that load packages and define macros
\input{includes} % various packages needed for maths etc.
\input{notation} % short-hand notation and macros


%%%%%%%%%%%%%%%%%%%%%%%%%%%%

\begin{document}
% front page



\input{titlepage}


%%%%%%%%%%%%%%%%%%%%%%%%%%% table of content
%If a table of content is needed, simply uncomment the following lines
%\tableofcontents
%\newpage

%%%%%%%%%%%%%%%%%%%%%%%%%%%% Main document
\section{Propuesta de Monografía}
[Este párrafo debe incluir un resumen conciso de la monografía, presentando el tema específico que se aborda. Describe su origen, antecedentes teóricos y prácticos.Se debe presenta la metodología que se utilizará para llevar a cabo el estudio, explicando el tipo de investigación y la técnica de análisis empleada. También se deben mencionar las teorías o referencias clave que respaldarán la investigación, incluyendo una cita a un trabajo específico, como \cite{Ejemplo}.

El párrafo también debe abordar las limitaciones de la investigación, como los obstáculos que pueden surgir y cómo se manejarán. Además, debe describir el alcance del estudio, especificando las áreas o aspectos específicos en los que se centrará la investigación. ]

\section{Identificación del Problema}
[Aquí se debe describir claramente el problema o la pregunta de investigación que se abordará en la monografía, por ejemplo, entender el desarrollo articulo y sus posibles aplicaciones, la solución a una pregunta o la aplicación de alguna demostración para generar algoritmos programables , es decir analizar las causas que dan importancia al desarrollo de este trabajo y a partir de allí, construir preguntas que guíen el proceso de
investigación. ]

\section{Objetivos de la Monografía}

[Aquí se deben establecer los objetivos generales y específicos de la monografía, indicando qué se espera lograr con la investigación, por ejemplo, crear un código de programación publico basado en algún algoritmo, solucionar una pregunta sobre el desarrollo de un articulo o describir una aplicación basada en modelos matemáticos.]

\section{Justificación}

[En esta sección se debe exponer la relevancia y la importancia de abordar el problema o la pregunta de investigación propuesta
argumentando que, no existe una solución para la problemática anunciada y que, por lo
tanto, se requiere de un abordaje especifico para la consecución de una solución.]






\section{Plan de Trabajo}

[Aquí se debe presentar un cronograma detallado con las actividades a realizar durante el desarrollo de la monografía, indicando las fechas de inicio y fin de cada actividad.
]
%(Al usar mas de una tabla se debe crear una referencia)

\begin{table}[htbp]
\centering
\begin{tabular}{@{}ll@{}}
\toprule
\textbf{Actividades} & \textbf{Estado} \\ 
\midrule
Planteamiento del problema & Finalizado \\ 
Revisión Bibliográfica & \\ 
Desarrollo del problema & \\ 
Redacción monografía & \\ 
Revisión  por docente Asesor & \\ 
Corrección de monografía & \\ 
Entrega del documento al docente evaluador & \\ 
Entrega del documento corregido al docente evaluador \hspace{2cm} & \\ 
Socialización del Trabajo de Grado & En proceso \\ \bottomrule
\end{tabular}
\caption{Estado de las actividades}
\label{tab:estado_actividades}
\end{table}





\section{Cronograma}

[Detallar los pasos que se seguirán para completar la monografía, incluyendo fechas tentativas de entrega de borradores, reuniones con el docente, y cualquier otro aspecto relevante del proceso.]




\begin{table}[htbp]
\centering
\resizebox{\textwidth}{!}{%
\begin{tabular}{|l|*{16}{c|}}
\hline
\multirow{2}{*}{
           Actividades
} & \multicolumn{16}{c|}{Semanas} \\ \cline{2-17} 
 & 1 & 2 & 3 & 4 & 5 & 6 & 7 & 8 & 9 & 10 & 11 & 12 & 13 & 14 & 15 & 16 \\ \hline
\multicolumn{1}{|c|}{Revision Bibliografica} &  &  &  &  &  &  &  &  &  &  &  &  &  &  &  &  \\ \hline
\multicolumn{1}{|c|}{Desarrollo del problema} &  &  &  &  &  &  &  &  &  &  &  &  &  &  &  &  \\ \hline
\multicolumn{1}{|c|}{Redaccion monografia} &  &  &  &  &  & &  &  &  &  &  &  &  &  &  &  \\ \hline
\multicolumn{1}{|c|}{Revisión por docente asesor} &  &  &  &  &  &  &  &  &  & &  &  &  &  &  &  \\ \hline
\multicolumn{1}{|c|}{Corrección de monografía} &  &  &  &  &  &  &  &  &  &  &  &  &  &  &  &  \\ \hline
\multicolumn{1}{|c|}{Revisión por el docente evaluador} &  &  &  &  &  &  &  &  &  &  &  & X &X  &  & &  \\ \hline
\multicolumn{1}{|c|}{Entrega del documento corregido al docente evaluador} &  &  &  &  &  &  &  &  &  &  &  &   & X &X  & X &  \\


\hline
\multicolumn{1}{|c|}{Socialización del Trabajo de Grado} &  &  &  &  &  &  &  &  &  &  &  &  &  &  & X & X \\ \hline
\end{tabular}%
}
\caption{Planificación de actividades}
\label{tab:planificacion}
\end{table}


\section*{Firmas}

\begin{flushleft}

Firma del Estudiante 1:
\hrulefill \hspace{1 cm} Fecha: \hrulefill \\

\vspace{1 cm}

Firma del Estudiante 2: \hrulefill \hspace{1 cm} Fecha: \hrulefill\\

\vspace{1 cm}

Firma del Docente: \hrulefill \hspace{1 cm} \vspace{1 cm}    Fecha: \hrulefill\\

\end{flushleft}


\renewcommand{\refname}{Referencias} % Cambia el título de la sección de bibliografía


\bibliography{bib} % Especifica el archivo .bib donde tienes tus referencias


\end{document}

%Final
